\documentclass[10pt,a4paper]{amsart}
\usepackage[margin=19mm]{geometry}
\usepackage{amsmath,amssymb,mathtools}
\usepackage[scr=rsfs,cal=euler]{mathalfa}
\usepackage{xfrac}
\usepackage[unicode,hidelinks,bookmarks=false]{hyperref}
\newcommand{\ie}{i.e.\ }
\newcommand{\cf}{cf.\ }
\newcommand{\resp}{respectively\ }
\newcommand{\Subsection}{\subsection}
\providecommand{\nobreakdash}{\nobreak\hbox{-}}
\setlength{\emergencystretch}{2em}
\multlinegap0pt
\usepackage{amssymb}
\usepackage{mathtools}
\usepackage[shortlabels]{enumitem}
\newtheorem{proposition}{Proposition}
\newtheorem{theorem}{Theorem}
\title{\textbf{Shape Optimization and Bernoulli Free Boundary Problems\\ for the Riemannian $p$-Laplacian}}
\author{Ababacar Sadikhe Djite, Diaraf Seck}
\date{Source: arXiv:2609.00571v2 (CC BY 4.0)}
\begin{document}
\maketitle

\noindent\emph{Source and licence.} \textbf{Shape Optimization and Bernoulli Free Boundary Problems\\ for the Riemannian $p$-Laplacian}. Version arXiv:2609.00571v2 (CC BY 4.0), distributed by arXiv under the Creative Commons Attribution 4.0 International licence, http://creativecommons.org/licenses/by/4.0/, as recorded in the arXiv licence metadata for this exact version and shown on the abstract page https://arxiv.org/abs/2609.00571v2. Full licence notice: \texttt{licenses/arxiv-2609.00571-CC-BY-4.0.txt}.\par\smallskip

\noindent\textit{Editorial scope.} Counted unit: the article's theorem thm:monotonicity-c-riemannian (source0.tex lines 2611--2630). The article's complete original proof of it is delivered with it (source0.tex lines 2632--2835, one \texttt{proof} environment of 204 lines). No mathematics is written, repaired or re-derived: the statement and the proof are the article's own lines, in the article's own notation, and no retained line is substituted or re-flowed.  Retained uncounted as the source-local context the proof needs: lines 1725--1764.  Nothing was omitted from any retained span, so the package carries no omission record.  No retained line contains a citation, so the wrapper carries no bibliography and no bibliography entry.  No retained line contains a figure environment, an image-inclusion command or drawing code, so no image is delivered and the package declares no asset dependency.  Editorial formatting only, in addition: an AMS \texttt{amsart} preamble that restates the article's own declarations and macros used by the retained lines, namely the environments \texttt{proposition}, \texttt{theorem} on separate counters, together with the standard packages those lines need.  The retained environments print in this document as Proposition 1, Theorem 1 in order of appearance, whereas the article prints the same items with the numbering of its own full document.  The complete native source of the article as distributed in the arXiv e-print for this exact version is bundled unchanged as \texttt{sources/209040/source0.tex}.
\begin{proposition}[Riemannian comparison at a contact point]
\label{prop:riemannian-contact-comparison}
Let
\[
K\subset\Omega_1\subset\Omega_2
\]
be admissible domains satisfying the Riemannian contraction hypothesis
{\rm (RC)}. Let
\[
u_i:=u_{\Omega_i},
\qquad i=1,2,
\]
and let
\[
v=u_2\circ\Psi^{-1}
\]
be the transported state in $\Omega_\Psi\setminus K$, where
$\Psi$ and the contact point $x_0$ are given by {\rm (RC)}. Then
\[
-\partial_{\nu_g}u_1(x_0)
>
-\partial_{\nu_g}v(x_0).
\label{eq:hopf-contact-comparison}
\]
Moreover, {\rm (RC)} gives
\[
-\partial_{\nu_g}v(x_0)
\geq
\kappa_\Psi
\bigl(-\partial_{\nu_g}u_2(y_0)\bigr),
\qquad
y_0=\Psi^{-1}(x_0),
\label{eq:RC-flux-comparison}
\]
with $\kappa_\Psi>1$. Consequently,
\begin{equation}
c_{\Omega_1}>\kappa_\Psi c_{\Omega_2}>c_{\Omega_2}.
\label{eq:riemannian-c-monotonicity}
\end{equation}
\end{proposition}

\begin{theorem}[Monotonicity of the Bernoulli constant]
\label{thm:monotonicity-c-riemannian}
Under the hypotheses of Proposition~\ref{prop:riemannian-contact-comparison},
if
\[
K\subset\Omega_1\subset\Omega_2,
\]
then
\[
c_{\Omega_1}>c_{\Omega_2}.
\]
Thus, along a decreasing sequence of admissible domains
\[
\Omega_{j+1}\subsetneq\Omega_j,
\]
the associated Bernoulli constants satisfy
\[
c_{j+1}<c_j.
\]
\end{theorem}

\begin{proof}
Set
\[
u_1:=u_{\Omega_1},
\qquad
u_2:=u_{\Omega_2}.
\]
By the Riemannian contraction hypothesis {\rm (RC)}, there exist an
admissible domain $\Omega_\Psi\subset\Omega_1$, a $C^2$ diffeomorphism
\[
\Psi:
\overline{\Omega_2\setminus K}
\longrightarrow
\overline{\Omega_\Psi\setminus K},
\]
and a contact point
\[
x_0\in\partial\Omega_\Psi\cap\partial\Omega_1
\]
such that
\[
\nu_{\Omega_\Psi}(x_0)=\nu_{\Omega_1}(x_0).
\]
Define the transported state
\[
v:=u_2\circ\Psi^{-1}
\qquad\text{in }\Omega_\Psi\setminus K.
\]
The hypothesis {\rm (RC)} asserts that $v$ is a weak subsolution of the
Riemannian $p$-Laplace equation and that
\[
v=1\quad\text{on }\partial K,
\qquad
v=0\quad\text{on }\partial\Omega_\Psi.
\]

\medskip
\noindent\textbf{Step 1: comparison of $v$ with $u_1$.}
Since $\Omega_\Psi\subset\Omega_1$, the function $u_1$ is
$p$-harmonic in $\Omega_\Psi\setminus K$ and satisfies
\[
u_1=1\quad\text{on }\partial K,
\qquad
u_1=0\quad\text{on }\partial\Omega_1.
\]
In particular,
\[
u_1\geq0\quad\text{on }\partial\Omega_\Psi,
\]
because $\partial\Omega_\Psi$ lies inside $\Omega_1$ except at the
contact point $x_0$, where $u_1(x_0)=0$. Hence
\[
(v-u_1)^+=0
\quad\text{on }\partial(\Omega_\Psi\setminus K).
\]
Thus
\[
w:=(v-u_1)^+\in W_0^{1,p}(\Omega_\Psi\setminus K).
\]

Using $w$ as an admissible test function in the weak subsolution
inequality for $v$ and in the weak equation for $u_1$, and subtracting
the two relations, gives
\[
\int_{\{v>u_1\}}
\left\langle
|\nabla_gv|^{p-2}\nabla_gv
-
|\nabla_gu_1|^{p-2}\nabla_gu_1,
\nabla_gv-\nabla_gu_1
\right\rangle_g\,dV_g
\leq0.
\]
The map
\[
\xi\longmapsto|\xi|_g^{p-2}\xi
\]
is strictly monotone for $1<p<\infty$. Therefore the integrand is
nonnegative and vanishes only when
\[
\nabla_gv=\nabla_gu_1.
\]
It follows that
\[
\nabla_g(v-u_1)^+=0
\quad\text{a.e. in }\Omega_\Psi\setminus K.
\]
Since $(v-u_1)^+$ has zero trace, we conclude
\[
v\leq u_1
\qquad\text{in }\Omega_\Psi\setminus K.
\]

\medskip
\noindent\textbf{Step 2: strict boundary comparison at the contact point.}
At the contact point $x_0$,
\[
v(x_0)=0=u_1(x_0).
\]
If $v\not\equiv u_1$, the strong comparison principle gives
\[
v<u_1
\qquad\text{in the interior of }\Omega_\Psi\setminus K.
\]
Both functions vanish at $x_0$ on the common tangent boundary, and
{\rm (RC)} gives
\[
\nu_{\Omega_\Psi}(x_0)=\nu_{\Omega_1}(x_0).
\]
The boundary point lemma of Hopf, applied with this common outward
normal, therefore yields
\[
-\partial_{\nu_g}u_1(x_0)
>
-\partial_{\nu_g}v(x_0).
\]
Since $\Omega_1$ is an optimal domain, its overdetermined boundary
condition is
\[
-\partial_{\nu_g}u_1=c_{\Omega_1}
\qquad\text{on }\partial\Omega_1.
\]
Consequently,
\[
c_{\Omega_1}>
-\partial_{\nu_g}v(x_0).
\]

\medskip
\noindent\textbf{Step 3: use of the flux amplification in {\rm (RC)}.}
Let
\[
y_0:=\Psi^{-1}(x_0).
\]
Because $\Omega_2$ is also an optimal domain,
\[
-\partial_{\nu_g}u_2(y_0)=c_{\Omega_2}.
\]
The flux part of {\rm (RC)} gives a constant $\kappa_\Psi>1$ such that
\[
-\partial_{\nu_g}v(x_0)
\geq
\kappa_\Psi
\left(
-\partial_{\nu_g}u_2(y_0)
\right)
=
\kappa_\Psi c_{\Omega_2}.
\]
Combining this inequality with the strict Hopf comparison obtained in
Step 2 gives the key chain
\[
c_{\Omega_1}
>
-\partial_{\nu_g}v(x_0)
\geq
\kappa_\Psi c_{\Omega_2}
>
c_{\Omega_2}.
\]
In particular,
\[
c_{\Omega_1}>c_{\Omega_2}.
\]

\medskip
\noindent\textbf{Step 4: return to the Lagrange multipliers.}
For every optimal domain,
\[
\lambda_\Omega
=
-\frac{p-1}{p}c_\Omega^p.
\]
The function
\[
c\longmapsto-\frac{p-1}{p}c^p
\]
is strictly decreasing on $(0,\infty)$. Hence
\[
c_{\Omega_1}>c_{\Omega_2}
\quad\Longrightarrow\quad
\lambda_{\Omega_1}<\lambda_{\Omega_2}.
\]
Thus the two formulations are equivalent:
\[
c_{\Omega_1}>c_{\Omega_2}
\qquad\Longleftrightarrow\qquad
\lambda_{\Omega_1}<\lambda_{\Omega_2}.
\]

Finally, applying the result to two consecutive domains of the
successive approximation,
\[
\Omega_{j+1}\subsetneq\Omega_j,
\]
gives
\[
c_{j+1}<c_j,
\qquad
\lambda_{j+1}>\lambda_j.
\]
This is precisely the monotonicity needed in the subsequent stopping
argument.
\end{proof}

\end{document}
